\chapter{模板使用说明与排版规范}

本章节旨在详细说明东北大学学位论文 LaTeX 模板的使用方法。本模板基于 `neudissertation.cls` 构建，预置了符合学校规范的字体、字号、行距及图表样式 。

\section{文档设置与学位适配}

本模板支持学士（Bachelor）、硕士（Master）和博士（Doctor）三种学位论文的排版。根据学位类型的不同，需要在文档类参数、个人信息变量以及文末章节进行相应的调整。

\subsection{文档类型参数 (Type)}

在 \texttt{main.tex} 的第一行，通过 \texttt{Type} 参数指定学位类型。该参数决定了封面样式、页眉页脚格式以及必要的元数据字段。

\begin{itemize}
    \item \textbf{本科论文}：\verb|\documentclass[Type=Bachelor]{neudissertation}|
    \item \textbf{硕士论文}：\verb|\documentclass[Type=Master]{neudissertation}|
    \item \textbf{博士论文}：\verb|\documentclass[Type=Doctor]{neudissertation}|
\end{itemize}

\subsection{必要变量设置}

所有学位类型的通用信息（如标题、作者、导师、学院等）在 \texttt{main.tex} 的“基础信息”和“导师”部分设置。

以下变量直接决定了中英文封面的内容，请务必准确填写：

\begin{table}[H]
  \centering
  \twotabcaption{tab:basic-vars}{基础信息变量表}{Basic Information Variables}
  \begin{tabular}{lll}
    \toprule
    类别 & 变量命令 & 说明 \\
    \midrule
    \multirow{2}{*}{标题} & \verb|\zhtitle{...}| & 中文论文题目 \\
    & \verb|\entitle{...}| & 英文论文题目 \\
    \midrule
    \multirow{2}{*}{作者} & \verb|\zhauthor{...}| & 作者中文姓名 \\
    & \verb|\enauthor{...}| & 作者英文姓名（拼音） \\
    \midrule
    \multirow{3}{*}{学籍} & \verb|\studentid{...}| & 学号 \\
    & \verb|\campusname{...}| & 学院名称（如：软件学院） \\
    & \verb|\zhmajor{...}| / \verb|\enmajor{...}| & 专业名称（中/英） \\
    \midrule
    时间 & \verb|\zhthesisdate{...}| & 封面日期（如：202x年6月） \\
    \bottomrule
  \end{tabular}
\end{table}

模板支持\textbf{单导师}和\textbf{双导师}模式。请注意，英文导师姓名通常需要包含职称（如 Professor X）。

\textbf{1. 第一导师（必须设置）}

\begin{verbatim}
\zhsupervisor{李四}                % 中文姓名
\zhsupervisortitle{教授}           % 职称
\zhsupervisorwork{东北大学软件学院} % 工作单位
\ensupervisor{Professor Li Si}     % 英文姓名(含职称)
\end{verbatim}

\textbf{2. 第二导师（可选设置）}

若您的论文有校外导师或协助指导教师，请填写 \verb|zhsupervisortwo| 系列变量。

\begin{itemize}
  \item \textbf{有第二导师时}：正常填写下列变量。
  \item \textbf{无第二导师时}：\textbf{必须}在 \texttt{main.tex} 中注释掉或删除下列四行代码！否则封面上会出现多余的空白或标点符号。
\end{itemize}

\begin{verbatim}
% --- 只有在有第二导师时才取消注释 ---
% \zhsupervisortwo{张三}
% \zhsupervisortwotitle{副教授}
% \zhsupervisortwowork{东北大学计算机学院}
% \ensupervisortwo{Associate Prof. Zhang San}
% ---------------------------------
\end{verbatim}

对于\textbf{硕士和博士}学位论文，还必须补充填写“仅硕博需要下列变量”区域的信息，否则封面和题名页的相关字段将为空白：

\begin{table}[H]
  \centering
  \twotabcaption{tab:grad-vars}{硕博学位论文专属变量}{Variables for Master and Doctoral Dissertations}
  \begin{tabular}{lll}
    \toprule
    变量命令 & 含义 & 示例 \\
    \midrule
    \verb|\zhdegreetype| & 学位类别 & 硕士 / 博士 \\
    \verb|\zhmajortype| & 学科门类 & 工学 / 理学 \\
    \verb|\zhsubmissiondate| & 论文提交日期 & 20xx年6月 \\
    \verb|\zhdefensedate| & 论文答辩日期 & 20xx年6月 \\
    \verb|\zhdegreegivendate| & 学位授予日期 & 20xx年6月 \\
    \verb|\zhdefensechairmen| & 答辩委员会主席 & 某某教授 \\
    \verb|\zhreviewers| & 评阅人 & 评阅人1\quad 评阅人2 \\
    \bottomrule
  \end{tabular}
\end{table}

封面左上角的分类号和密级一般留空即可：

\begin{itemize}
    \item \verb|\classification{}|：留空。
    \item \verb|\serialno{}|：留空。
    \item \verb|\udc{}|：留空。
\end{itemize}

\subsection{硕博补充章节：科研及发表论文}

根据学校要求，硕士和博士论文需要在参考文献之后、致谢之前，增加“攻读学位期间科研及发表论文情况”章节。

请在 \texttt{main.tex} 的 \verb|\backmatter| 区域（参考文献之后），取消相关代码的注释：

\begin{verbatim}
% --------------------------------------------------
% 仅硕博需要此章节
\chapter{攻读学位期间科研及发表论文情况}
{
   \noindent \sanhao \bf 科研情况:
}
\begin{enumerate}[label*=项目\arabic*：, leftmargin=*,
    labelindent=\parindent, itemindent=6em]
    \item 国家自然科学基金项目...
\end{enumerate}

{
    \noindent \sanhao \bf 论文发表情况:
}

\begin{enumerate}[label*=发表\arabic*：,  leftmargin=*,
    labelindent=\parindent,  itemindent=6em]
    \item 作者名. 论文题目...
\end{enumerate}
% -------------------------------------------------
\end{verbatim}

\textbf{注意}：本科毕业设计（论文）不需要此章节，请保持注释状态。

\section{字体与文本设置}

\subsection{中文字体}
模板预定义了符合规范的中文字体族，通常情况下系统会自动根据章节标题（黑体）和正文（宋体）切换，但在需要强调时可手动指定 ：

\begin{table}[H]
  \centering
  \twotabcaption{tab:font-demo}{常用字体命令对照表}{Comparison of Common Font Commands}
  \begin{tabular}{llll}
    \toprule
    命令 & 说明 & 示例 & 备注 \\
    \midrule
    \verb|\song| & 宋体 & \song 东北大学 (Northeastern University) & 正文默认 \\
    \verb|\hei| & 黑体 & \hei 东北大学 (Northeastern University) & 标题默认 \\
    \verb|\kai| & 楷体 & \kai 东北大学 (Northeastern University) & 摘要/强调 \\
    \verb|\fs| & 仿宋 & \fs 东北大学 (Northeastern University) & 很少使用 \\
    \bottomrule
  \end{tabular}
\end{table}

\subsection{文字强调}
正文中请使用以下标准命令进行强调：
\begin{itemize}
    \item \textbf{加粗}：使用 \verb|\textbf{文字}|。
    \item \textit{斜体}：使用 \verb|\textit{文字}|（中文环境通常显示为楷体）。
\end{itemize}

\section{公式排版}

\subsection{公式环境}
行内公式请使用 \verb|$...$| 包裹，例如 $E=mc^2$。
行间公式推荐使用 \verb|equation| 环境，以便自动生成编号：

\begin{equation}
  f(x) = \frac{1}{\sqrt{2\pi}\sigma} \exp\left(-\frac{(x-\mu)^2}{2\sigma^2}\right)
  \label{eq:guassian}
\end{equation}

\subsection{公式引用}
引用公式时，请务必使用模板定义的 \verb|\equref{label}| 命令。
\begin{itemize}
  \item \textbf{正确示例}：如 \equref{eq:guassian} 所示（会自动添加“式”字前缀）。
  \item \textbf{错误示例}：如 公式\ref{eq:guassian} 所示。
\end{itemize}

\section{图表双语标题规范}

根据学校学位论文要求，图和表必须包含\textbf{中英双语标题}。本模板封装了 twofigcaption和 twotabcaption 命令来简化此操作 。

\subsection{图片的插入}
请注意：图片标题应位于图片\textbf{下方}。使用 \verb|\twofigcaption{label}{中文}{英文}| 命令。

\begin{figure}[H]
  \centering
  % 使用 demo 占位符，实际请替换为 \includegraphics
  \fbox{
    \begin{minipage}{0.6\textwidth}
      \centering
      \vspace{1cm}
      此处插入图片文件 \\
      Image Placeholder
      \vspace{1cm}
  \end{minipage}}

  \twofigcaption{fig:structure}{系统整体架构图}{System Architecture Diagram}
\end{figure}

引用图片请使用 \verb|\figref{fig:structure}|，效果为：\figref{fig:structure}。

有能力的建议使用tikz 绘图，可以获得媲美svg的显示效果，示例如下.

\begin{figure}[H]
  \centering
  \begin{tikzpicture}[
    every node/.style = {anchor=south west, align=center}
]
% \draw [help lines] (0,0) grid (14,5);
\node[name=train_model, draw, fill=medium blue] {模型训练};
\node[name=optimization, draw, fill=color2, right=of train_model] {轻量化};
\node[name=optimized_model, draw, fill=medium blue, right=3cm of optimization] {(高能效)\\模型推理};
\draw[dash dot] ($(train_model.south west)+(-0.2cm, -0.2cm)$)  |- ($(optimization.north east)+(0.2cm, 0.2cm)$) 
    node[pos=0.75, anchor=south] {训练环境} |- ($(train_model.south west)+(-0.2cm, -0.2cm)$);
\draw[dash dot] ($(optimized_model.south west)+(-0.2cm, -0.2cm)$) |- ($(optimized_model.north east)+(0.2cm, 0.2cm)$) node[pos=0.75, anchor=south]{移动环境} |- ($(optimized_model.south west)+(-0.2cm, -0.2cm)$);
\draw[arrow]  ($(optimization.east)+(0cm, 0.2cm)$) to[out=30, in=150] node[above]{模型部署} ($(optimized_model.west)+(0cm, 0.2cm)$);
\draw[arrow, dashed]  ($(optimized_model.west)+(0cm, -0.2cm)$) to[out=210, in=330] node[below]{迭代更新} ($(optimization.east)+(0cm, -0.2cm)$);
\draw[arrow] (train_model) -- (optimization);
\end{tikzpicture}
  \twofigcaption{fig:optimization-process}{深度神经网络轻量化与部署过程}{Lightweight and Deploy Process of Deep Neural Network}
\end{figure}

多图并列排版示例如下

\begin{figure}[H]
  \centering
  \begin{minipage}[t]{0.45\linewidth}
    \centering
    \resizebox{\textwidth}{!}{% Created by tikzDevice version 0.12.3.1 on 2022-03-02 20:32:42
% !TEX encoding = UTF-8 Unicode
\begin{tikzpicture}[x=1pt,y=1pt]
\definecolor{fillColor}{RGB}{255,255,255}
\path[use as bounding box,fill=fillColor,fill opacity=0.00] (0,0) rectangle (216.81,216.81);
\begin{scope}
\path[clip] (  0.00,  0.00) rectangle (216.81,216.81);
\definecolor{drawColor}{RGB}{255,255,255}
\definecolor{fillColor}{RGB}{255,255,255}

\path[draw=drawColor,line width= 0.6pt,line join=round,line cap=round,fill=fillColor] (  0.00,  0.00) rectangle (216.81,216.81);
\end{scope}
\begin{scope}
\path[clip] ( 36.11, 59.22) rectangle (211.31,174.86);
\definecolor{drawColor}{RGB}{0,0,0}
\definecolor{fillColor}{RGB}{255,255,255}

\path[draw=drawColor,line width= 0.6pt,line join=round,line cap=round,fill=fillColor] ( 36.11, 59.22) rectangle (211.31,174.86);
\definecolor{drawColor}{RGB}{255,255,255}

\path[draw=drawColor,line width= 0.3pt,line join=round] ( 36.11, 81.59) --
	(211.31, 81.59);

\path[draw=drawColor,line width= 0.3pt,line join=round] ( 36.11,115.84) --
	(211.31,115.84);

\path[draw=drawColor,line width= 0.3pt,line join=round] ( 36.11,150.08) --
	(211.31,150.08);

\path[draw=drawColor,line width= 0.6pt,line join=round] ( 36.11, 64.47) --
	(211.31, 64.47);

\path[draw=drawColor,line width= 0.6pt,line join=round] ( 36.11, 98.72) --
	(211.31, 98.72);

\path[draw=drawColor,line width= 0.6pt,line join=round] ( 36.11,132.96) --
	(211.31,132.96);

\path[draw=drawColor,line width= 0.6pt,line join=round] ( 36.11,167.20) --
	(211.31,167.20);

\path[draw=drawColor,line width= 0.6pt,line join=round] ( 47.54, 59.22) --
	( 47.54,174.86);

\path[draw=drawColor,line width= 0.6pt,line join=round] ( 66.58, 59.22) --
	( 66.58,174.86);

\path[draw=drawColor,line width= 0.6pt,line join=round] ( 85.62, 59.22) --
	( 85.62,174.86);

\path[draw=drawColor,line width= 0.6pt,line join=round] (104.67, 59.22) --
	(104.67,174.86);

\path[draw=drawColor,line width= 0.6pt,line join=round] (123.71, 59.22) --
	(123.71,174.86);

\path[draw=drawColor,line width= 0.6pt,line join=round] (142.75, 59.22) --
	(142.75,174.86);

\path[draw=drawColor,line width= 0.6pt,line join=round] (161.80, 59.22) --
	(161.80,174.86);

\path[draw=drawColor,line width= 0.6pt,line join=round] (180.84, 59.22) --
	(180.84,174.86);

\path[draw=drawColor,line width= 0.6pt,line join=round] (199.88, 59.22) --
	(199.88,174.86);
\definecolor{fillColor}{RGB}{158,214,229}

\path[fill=fillColor] ( 42.78, 64.47) rectangle ( 52.30, 66.80);

\path[fill=fillColor] ( 61.82, 64.47) rectangle ( 71.34, 72.35);

\path[fill=fillColor] ( 80.86, 64.47) rectangle ( 90.38, 77.83);

\path[fill=fillColor] ( 99.91, 64.47) rectangle (109.43, 87.07);
\definecolor{fillColor}{RGB}{252,157,154}

\path[fill=fillColor] (118.95, 64.47) rectangle (128.47, 93.92);

\path[fill=fillColor] (137.99, 64.47) rectangle (147.51,169.60);

\path[fill=fillColor] (157.04, 64.47) rectangle (166.56, 93.92);

\path[fill=fillColor] (176.08, 64.47) rectangle (185.60, 94.61);

\path[fill=fillColor] (195.12, 64.47) rectangle (204.64,131.93);
\end{scope}
\begin{scope}
\path[clip] (  0.00,  0.00) rectangle (216.81,216.81);
\definecolor{drawColor}{gray}{0.30}

\node[text=drawColor,anchor=base east,inner sep=0pt, outer sep=0pt, scale=  0.88] at ( 31.16, 61.44) {0};

\node[text=drawColor,anchor=base east,inner sep=0pt, outer sep=0pt, scale=  0.88] at ( 31.16, 95.68) {100};

\node[text=drawColor,anchor=base east,inner sep=0pt, outer sep=0pt, scale=  0.88] at ( 31.16,129.93) {200};

\node[text=drawColor,anchor=base east,inner sep=0pt, outer sep=0pt, scale=  0.88] at ( 31.16,164.17) {300};
\end{scope}
\begin{scope}
\path[clip] (  0.00,  0.00) rectangle (216.81,216.81);
\definecolor{drawColor}{gray}{0.20}

\path[draw=drawColor,line width= 0.6pt,line join=round] ( 33.36, 64.47) --
	( 36.11, 64.47);

\path[draw=drawColor,line width= 0.6pt,line join=round] ( 33.36, 98.72) --
	( 36.11, 98.72);

\path[draw=drawColor,line width= 0.6pt,line join=round] ( 33.36,132.96) --
	( 36.11,132.96);

\path[draw=drawColor,line width= 0.6pt,line join=round] ( 33.36,167.20) --
	( 36.11,167.20);
\end{scope}
\begin{scope}
\path[clip] (  0.00,  0.00) rectangle (216.81,216.81);
\definecolor{drawColor}{gray}{0.20}

\path[draw=drawColor,line width= 0.6pt,line join=round] ( 47.54, 56.47) --
	( 47.54, 59.22);

\path[draw=drawColor,line width= 0.6pt,line join=round] ( 66.58, 56.47) --
	( 66.58, 59.22);

\path[draw=drawColor,line width= 0.6pt,line join=round] ( 85.62, 56.47) --
	( 85.62, 59.22);

\path[draw=drawColor,line width= 0.6pt,line join=round] (104.67, 56.47) --
	(104.67, 59.22);

\path[draw=drawColor,line width= 0.6pt,line join=round] (123.71, 56.47) --
	(123.71, 59.22);

\path[draw=drawColor,line width= 0.6pt,line join=round] (142.75, 56.47) --
	(142.75, 59.22);

\path[draw=drawColor,line width= 0.6pt,line join=round] (161.80, 56.47) --
	(161.80, 59.22);

\path[draw=drawColor,line width= 0.6pt,line join=round] (180.84, 56.47) --
	(180.84, 59.22);

\path[draw=drawColor,line width= 0.6pt,line join=round] (199.88, 56.47) --
	(199.88, 59.22);
\end{scope}
\begin{scope}
\path[clip] (  0.00,  0.00) rectangle (216.81,216.81);
\definecolor{drawColor}{gray}{0.30}

\node[text=drawColor,rotate= 30.00,anchor=base east,inner sep=0pt, outer sep=0pt, scale=  0.88] at ( 50.57, 49.02) {GoogleNet};

\node[text=drawColor,rotate= 30.00,anchor=base east,inner sep=0pt, outer sep=0pt, scale=  0.88] at ( 69.61, 49.02) {ResNet50};

\node[text=drawColor,rotate= 30.00,anchor=base east,inner sep=0pt, outer sep=0pt, scale=  0.88] at ( 88.65, 49.02) {RegNetY-8G};

\node[text=drawColor,rotate= 30.00,anchor=base east,inner sep=0pt, outer sep=0pt, scale=  0.88] at (107.70, 49.02) {EfficientNet-B7};

\node[text=drawColor,rotate= 30.00,anchor=base east,inner sep=0pt, outer sep=0pt, scale=  0.88] at (126.74, 49.02) {ViT-B};

\node[text=drawColor,rotate= 30.00,anchor=base east,inner sep=0pt, outer sep=0pt, scale=  0.88] at (145.78, 49.02) {ViT-L};

\node[text=drawColor,rotate= 30.00,anchor=base east,inner sep=0pt, outer sep=0pt, scale=  0.88] at (164.83, 49.02) {DeiT-B};

\node[text=drawColor,rotate= 30.00,anchor=base east,inner sep=0pt, outer sep=0pt, scale=  0.88] at (183.87, 49.02) {Swin-B};

\node[text=drawColor,rotate= 30.00,anchor=base east,inner sep=0pt, outer sep=0pt, scale=  0.88] at (202.91, 49.02) {Swin-L};
\end{scope}
\begin{scope}
\path[clip] (  0.00,  0.00) rectangle (216.81,216.81);
\definecolor{drawColor}{RGB}{0,0,0}

\node[text=drawColor,anchor=base,inner sep=0pt, outer sep=0pt, scale=  1.10] at (123.71,  7.64) {计算机视觉神经网络};
\end{scope}
\begin{scope}
%\path[clip] (  0.00,  0.00) rectangle (216.81,216.81);
\definecolor{drawColor}{RGB}{0,0,0}

\node[text=drawColor,rotate= 90.00,anchor=base, yshift=0.3cm,inner sep=0pt, outer sep=0pt, scale=  1.10] at ( 13.08,117.04) {参数量(百万)};
\end{scope}
\begin{scope}
\path[clip] (  0.00,  0.00) rectangle (216.81,216.81);
\definecolor{fillColor}{RGB}{255,255,255}

\path[fill=fillColor] ( 62.71,185.86) rectangle (184.71,211.31);
\end{scope}
\begin{scope}
\path[clip] (  0.00,  0.00) rectangle (216.81,216.81);

\path[] ( 73.71,191.36) rectangle ( 88.16,205.81);
\end{scope}
\begin{scope}
\path[clip] (  0.00,  0.00) rectangle (216.81,216.81);
\definecolor{fillColor}{RGB}{158,214,229}

\path[fill=fillColor] ( 74.42,192.07) rectangle ( 87.45,205.10);
\end{scope}
\begin{scope}
\path[clip] (  0.00,  0.00) rectangle (216.81,216.81);

\path[] (118.71,191.36) rectangle (133.17,205.81);
\end{scope}
\begin{scope}
\path[clip] (  0.00,  0.00) rectangle (216.81,216.81);
\definecolor{fillColor}{RGB}{252,157,154}

\path[fill=fillColor] (119.43,192.07) rectangle (132.46,205.10);
\end{scope}
\begin{scope}
\path[clip] (  0.00,  0.00) rectangle (216.81,216.81);
\definecolor{drawColor}{RGB}{0,0,0}

\node[text=drawColor,anchor=base west,inner sep=0pt, outer sep=0pt, scale=  0.88] at ( 93.66,195.55) {CNN};
\end{scope}
\begin{scope}
\path[clip] (  0.00,  0.00) rectangle (216.81,216.81);
\definecolor{drawColor}{RGB}{0,0,0}

\node[text=drawColor,anchor=base west,inner sep=0pt, outer sep=0pt, scale=  0.88] at (138.67,195.55) {X-formers};
\end{scope}
\end{tikzpicture}
}
  \end{minipage}%
  ~
  \begin{minipage}[t]{0.45\linewidth}
    \centering
    \resizebox{\textwidth}{!}{% Created by tikzDevice version 0.12.3.1 on 2022-03-01 01:30:14
% !TEX encoding = UTF-8 Unicode
\begin{tikzpicture}[x=1pt,y=1pt]
\definecolor{fillColor}{RGB}{255,255,255}
\path[use as bounding box,fill=fillColor,fill opacity=0.00] (0,0) rectangle (216.81,216.81);
\begin{scope}
\path[clip] (  0.00,  0.00) rectangle (216.81,216.81);
\definecolor{drawColor}{RGB}{255,255,255}
\definecolor{fillColor}{RGB}{255,255,255}

\path[draw=drawColor,line width= 0.6pt,line join=round,line cap=round,fill=fillColor] (  0.00,  0.00) rectangle (216.81,216.81);
\end{scope}
\begin{scope}
\path[clip] ( 46.86, 56.75) rectangle (211.31,174.86);
\definecolor{drawColor}{RGB}{0,0,0}
\definecolor{fillColor}{RGB}{255,255,255}

\path[draw=drawColor,line width= 0.6pt,line join=round,line cap=round,fill=fillColor] ( 46.86, 56.75) rectangle (211.31,174.86);
\definecolor{drawColor}{RGB}{255,255,255}

\path[draw=drawColor,line width= 0.3pt,line join=round] ( 46.86, 62.12) --
	(211.31, 62.12);

\path[draw=drawColor,line width= 0.3pt,line join=round] ( 46.86,103.07) --
	(211.31,103.07);

\path[draw=drawColor,line width= 0.3pt,line join=round] ( 46.86,144.03) --
	(211.31,144.03);

\path[draw=drawColor,line width= 0.6pt,line join=round] ( 46.86, 82.60) --
	(211.31, 82.60);

\path[draw=drawColor,line width= 0.6pt,line join=round] ( 46.86,123.55) --
	(211.31,123.55);

\path[draw=drawColor,line width= 0.6pt,line join=round] ( 46.86,164.51) --
	(211.31,164.51);

\path[draw=drawColor,line width= 0.6pt,line join=round] ( 58.90, 56.75) --
	( 58.90,174.86);

\path[draw=drawColor,line width= 0.6pt,line join=round] ( 78.95, 56.75) --
	( 78.95,174.86);

\path[draw=drawColor,line width= 0.6pt,line join=round] ( 99.01, 56.75) --
	( 99.01,174.86);

\path[draw=drawColor,line width= 0.6pt,line join=round] (119.06, 56.75) --
	(119.06,174.86);

\path[draw=drawColor,line width= 0.6pt,line join=round] (139.11, 56.75) --
	(139.11,174.86);

\path[draw=drawColor,line width= 0.6pt,line join=round] (159.17, 56.75) --
	(159.17,174.86);

\path[draw=drawColor,line width= 0.6pt,line join=round] (179.22, 56.75) --
	(179.22,174.86);

\path[draw=drawColor,line width= 0.6pt,line join=round] (199.28, 56.75) --
	(199.28,174.86);
\definecolor{fillColor}{RGB}{158,214,229}

\path[fill=fillColor] ( 53.88, 62.12) rectangle ( 63.91, 73.51);

\path[fill=fillColor] ( 73.94, 62.12) rectangle ( 83.96,117.39);
\definecolor{fillColor}{RGB}{252,157,154}

\path[fill=fillColor] ( 93.99, 62.12) rectangle (104.02,103.92);

\path[fill=fillColor] (114.05, 62.12) rectangle (124.07,113.96);

\path[fill=fillColor] (134.10, 62.12) rectangle (144.13,121.11);

\path[fill=fillColor] (154.16, 62.12) rectangle (164.18,125.89);

\path[fill=fillColor] (174.21, 62.12) rectangle (184.24,140.47);

\path[fill=fillColor] (194.26, 62.12) rectangle (204.29,169.49);
\end{scope}
\begin{scope}
\path[clip] (  0.00,  0.00) rectangle (216.81,216.81);
\definecolor{drawColor}{gray}{0.30}

\node[text=drawColor,anchor=base east,inner sep=0pt, outer sep=0pt, scale=  0.88] at ( 41.91, 79.57) {$10^{~}$};

\node[text=drawColor,anchor=base east,inner sep=0pt, outer sep=0pt, scale=  0.88] at ( 41.91,120.52) {$10^{3}$};

\node[text=drawColor,anchor=base east,inner sep=0pt, outer sep=0pt, scale=  0.88] at ( 41.91,161.48) {$10^{5}$};
\end{scope}
\begin{scope}
\path[clip] (  0.00,  0.00) rectangle (216.81,216.81);
\definecolor{drawColor}{gray}{0.20}

\path[draw=drawColor,line width= 0.6pt,line join=round] ( 44.11, 82.60) --
	( 46.86, 82.60);

\path[draw=drawColor,line width= 0.6pt,line join=round] ( 44.11,123.55) --
	( 46.86,123.55);

\path[draw=drawColor,line width= 0.6pt,line join=round] ( 44.11,164.51) --
	( 46.86,164.51);
\end{scope}
\begin{scope}
\path[clip] (  0.00,  0.00) rectangle (216.81,216.81);
\definecolor{drawColor}{gray}{0.20}

\path[draw=drawColor,line width= 0.6pt,line join=round] ( 58.90, 54.00) --
	( 58.90, 56.75);

\path[draw=drawColor,line width= 0.6pt,line join=round] ( 78.95, 54.00) --
	( 78.95, 56.75);

\path[draw=drawColor,line width= 0.6pt,line join=round] ( 99.01, 54.00) --
	( 99.01, 56.75);

\path[draw=drawColor,line width= 0.6pt,line join=round] (119.06, 54.00) --
	(119.06, 56.75);

\path[draw=drawColor,line width= 0.6pt,line join=round] (139.11, 54.00) --
	(139.11, 56.75);

\path[draw=drawColor,line width= 0.6pt,line join=round] (159.17, 54.00) --
	(159.17, 56.75);

\path[draw=drawColor,line width= 0.6pt,line join=round] (179.22, 54.00) --
	(179.22, 56.75);

\path[draw=drawColor,line width= 0.6pt,line join=round] (199.28, 54.00) --
	(199.28, 56.75);
\end{scope}
\begin{scope}
\path[clip] (  0.00,  0.00) rectangle (216.81,216.81);
\definecolor{drawColor}{gray}{0.30}

\node[text=drawColor,rotate= 30.00,anchor=base east,inner sep=0pt, outer sep=0pt, scale=  0.80] at ( 61.65, 47.03) {Cove-MTLSTM};

\node[text=drawColor,rotate= 30.00,anchor=base east,inner sep=0pt, outer sep=0pt, scale=  0.80] at ( 81.71, 47.03) {ELMo};

\node[text=drawColor,rotate= 30.00,anchor=base east,inner sep=0pt, outer sep=0pt, scale=  0.80] at (101.76, 47.03) {BERT-B};

\node[text=drawColor,rotate= 30.00,anchor=base east,inner sep=0pt, outer sep=0pt, scale=  0.80] at (121.81, 47.03) {BERT-L};

\node[text=drawColor,rotate= 30.00,anchor=base east,inner sep=0pt, outer sep=0pt, scale=  0.80] at (141.87, 47.03) {GPT3-L};

\node[text=drawColor,rotate= 30.00,anchor=base east,inner sep=0pt, outer sep=0pt, scale=  0.80] at (161.92, 47.03) {GPT3-XL};

\node[text=drawColor,rotate= 30.00,anchor=base east,inner sep=0pt, outer sep=0pt, scale=  0.80] at (181.98, 47.03) {GPT3-6.7B};

\node[text=drawColor,rotate= 30.00,anchor=base east,inner sep=0pt, outer sep=0pt, scale=  0.80] at (202.03, 47.03) {GPT3-170B};
\end{scope}
\begin{scope}
\path[clip] (  0.00,  0.00) rectangle (216.81,216.81);
\definecolor{drawColor}{RGB}{0,0,0}

\node[text=drawColor,anchor=base,inner sep=0pt, outer sep=0pt, scale=  1.10] at (129.09,  7.64) {自然语言处理神经网络};
\end{scope}
% \begin{scope}
% \path[clip] (  0.00,  0.00) rectangle (216.81,216.81);
% \definecolor{drawColor}{RGB}{0,0,0}

% \node[text=drawColor,rotate= 90.00,anchor=base,inner sep=0pt, outer sep=0pt, scale=  1.10] at ( 13.08,115.80) {数据量(百万)};
% \end{scope}
\begin{scope}
\path[clip] (  0.00,  0.00) rectangle (216.81,216.81);
\definecolor{fillColor}{RGB}{255,255,255}

\path[fill=fillColor] ( 68.02,185.86) rectangle (190.15,211.31);
\end{scope}
\begin{scope}
\path[clip] (  0.00,  0.00) rectangle (216.81,216.81);
\definecolor{fillColor}{gray}{0.95}

\path[fill=fillColor] ( 79.02,191.36) rectangle ( 93.48,205.81);
\end{scope}
\begin{scope}
\path[clip] (  0.00,  0.00) rectangle (216.81,216.81);
\definecolor{fillColor}{RGB}{158,214,229}

\path[fill=fillColor] ( 79.74,192.07) rectangle ( 92.77,205.10);
\end{scope}
\begin{scope}
\path[clip] (  0.00,  0.00) rectangle (216.81,216.81);
\definecolor{fillColor}{gray}{0.95}

\path[fill=fillColor] (124.15,191.36) rectangle (138.61,205.81);
\end{scope}
\begin{scope}
\path[clip] (  0.00,  0.00) rectangle (216.81,216.81);
\definecolor{fillColor}{RGB}{252,157,154}

\path[fill=fillColor] (124.86,192.07) rectangle (137.89,205.10);
\end{scope}
\begin{scope}
\path[clip] (  0.00,  0.00) rectangle (216.81,216.81);
\definecolor{drawColor}{RGB}{0,0,0}

\node[text=drawColor,anchor=base west,inner sep=0pt, outer sep=0pt, scale=  0.88] at ( 98.98,195.55) {RNN};
\end{scope}
\begin{scope}
\path[clip] (  0.00,  0.00) rectangle (216.81,216.81);
\definecolor{drawColor}{RGB}{0,0,0}

\node[text=drawColor,anchor=base west,inner sep=0pt, outer sep=0pt, scale=  0.88] at (144.11,195.55) {X-formers};
\end{scope}
\end{tikzpicture}

    }
  \end{minipage}
  \twofigcaption{fig:param-compare}{神经网络参数量对比}{Comparesion of Neural Network Parameter Amount}
\end{figure}

tikz绘图建议：Tikz绘图涉及多次调整。建议先在figure-pdf.tex中单独绘制。绘制完成后，创建一个tex文件，引入到主文档中。figure-pdf.tex和main.tex都共享tikzbase.sty。若要加入新的tikz 库引用，建议改动该文件，从而使tikz依赖一致。

\subsection{表格的插入}
请注意：表格标题应位于表格\textbf{上方}。使用 \verb|\twotabcaption{label}{中文}{英文}| 命令。推荐使用三线表（booktabs）。

\begin{table}[H]
  \centering
  \twotabcaption{tab:experiment}{实验数据对比}{Comparison of Experimental Data}
  \begin{tabular}{l c c c}
    \toprule
    模型 (Model) & 准确率 (Acc) & 召回率 (Recall) & F1值 \\
    \midrule
    Baseline & 85.2\% & 84.1\% & 0.846 \\
    Our Method & \textbf{92.1\%} & \textbf{91.5\%} & \textbf{0.918} \\
    \bottomrule
  \end{tabular}
\end{table}

引用表格请使用 \verb|\tabref{tab:experiment}|，效果为：\tabref{tab:experiment}。

\section{参考文献引用}

模板支持两种引用模式，请根据语境选择：

\begin{enumerate}
    \item \textbf{上标模式}（\verb|\upcite|）：
    \begin{itemize}
      \item 用于陈述事实或观点时。
      \item 示例：已有研究表明深度学习在图像识别中表现优异\upcite{guo_2023_research}。
    \end{itemize}

    \item \textbf{正文模式}（\verb|\cite|）：
    \begin{itemize}
      \item 用于文献作为句子成分（如主语）时。
      \item 示例：文献 \cite{guo_2023_research} 详细讨论了这一问题。
    \end{itemize}
\end{enumerate}

\section{算法伪代码}

对于计算机类专业的论文，常需插入算法伪代码。模板已预置 `algorithm` 和 `algorithmic` 宏包。

\begin{algorithm}[H]
  \caption{自适应学习率调整算法}
  \label{alg:learning_rate}
  \begin{algorithmic}[1]
    \REQUIRE Initial learning rate $\alpha_0$, decay rate $\gamma$
    \ENSURE Updated parameters $\theta$
    \STATE Initialize $t \leftarrow 0$
    \WHILE{not converged}
    \STATE $t \leftarrow t + 1$
    \STATE $\alpha_t \leftarrow \alpha_0 / (1 + \gamma t)$
    \STATE Update $\theta_t$ using gradient descent with $\alpha_t$
    \ENDWHILE
    \RETURN $\theta_t$
  \end{algorithmic}
\end{algorithm}

\section{其他注意事项}

\begin{itemize}
  \item \textbf{列表环境}：模板优化了 `itemize` 和 `enumerate` 的间距，使其更符合中文排版习惯 。
  \item \textbf{编译方式}：推荐使用Latex 结合-xelatex编译。确保中文字体和文献正确渲染。
\end{itemize}